\chapter{Planck pair, universal tension, and gravitational compliance}
\label{chap:pareja-planck-respuesta}

\section{Planck dimensional pair \((Q,\ell)\): \(Q\ell=\hbar c\), \(\ell/Q=G/c^4\), and unique reconstruction}

Let \(t_0\) be the chronological unit of the CT108 section, and let \(c\) be the
causal section selected by the constitutive reader. We define

\[
t_P=\frac{54}{\pi}t_0,
\qquad
\ell=\ell_P=ct_P=\frac{54}{\pi}ct_0,
\]

and the energy quantum

\[
Q=E_P=\frac{\hbar}{t_P}=\frac{h}{108t_0}.
\]

The periodicity \(T_{108}=108t_0\) verifies

\[
T_{108}=2\pi t_P,
\qquad
QT_{108}=h.
\]

From these definitions, the first identity of the pair is deduced:

\begin{equation}
\boxed{Q\ell=\hbar c.}
\label{eq:rm-planck-product}
\end{equation}

The circular gravitational realization determines

\[
G=\left(\frac{54}{\pi}\right)^2\frac{c^5t_0^2}{\hbar}.
\]

Substituting \(\ell=(54/\pi)ct_0\) and \(Q=\hbar/t_P\) yields the
second identity:

\begin{equation}
\boxed{\frac{\ell}{Q}=\frac{G}{c^4}.}
\label{eq:rm-planck-quotient}
\end{equation}

The equations \eqref{eq:rm-planck-product} and
\eqref{eq:rm-planck-quotient} are algebraically independent as
product and quotient relations. They uniquely reconstruct

\[
Q=\sqrt{\frac{\hbar c^5}{G}},
\qquad
\ell=\sqrt{\frac{G\hbar}{c^3}}.
\]

\begin{theorem}[Planck Pair Closure]
\label{thm:rm-planck-pair}
In the circular CT108 realization, the positive pair \((Q,\ell)\) is the
unique solution of
\[
Q\ell=\hbar c,
\qquad
\ell/Q=G/c^4.
\]
Satisfies, moreover,
\[
GQ=c^4\ell,
\qquad
GQ^2=\hbar c^5,
\qquad
G\hbar=c^3\ell^2.
\]
\end{theorem}

\begin{proof}
Positivity permits multiplication and division of the two equations without
ambiguity of sign. The first additional equality follows from
\(\ell/Q=G/c^4\); the second results from multiplying it by \(Q^2\), and the
third from multiplying \(\ell/Q=G/c^4\) by
\(Q\ell=\hbar c\). Reconstruction by positive roots proves
uniqueness.
\end{proof}

\section{Reciprocity between Planck tension and compliance and conservation of the action product}

The pair defines two reciprocal responses:

\begin{equation}
\mathscr T_P:=\frac{Q}{\ell}=\frac{c^4}{G},
\qquad
\mathscr C_P:=\frac{\ell}{Q}=\frac{G}{c^4}.
\label{eq:rm-tension-compliance}
\end{equation}

The first has the dimension of force; the second, the dimension of
gravitational compliance. They satisfy

\[
\mathscr T_P\mathscr C_P=1.
\]

The Einstein equation is then written as a constitutive law:

\begin{equation}
\boxed{
G_{\mu\nu}+\Lambda g_{\mu\nu}
=8\pi\mathscr C_P T_{\mu\nu}.}
\label{eq:rm-einstein-compliance}
\end{equation}

Matter supplies tension--energy, and geometry responds through
\(\mathscr C_P\). The constant \(G\) publishes the curvature capacity per
unit of tension in the chosen causal chart.

In the Newtonian regime, for an energy density
\(\varepsilon\),

\[
\nabla^2\!\left(\frac{\Phi}{c^2}\right)
=4\pi\mathscr C_P\varepsilon.
\]

The Einstein--Hilbert action is normalized by the structural area:

\begin{equation}
\frac{S_{\rm EH}}{\hbar}
=\frac{1}{16\pi\ell^2}
\int_{\mathcal M}(R-2\Lambda)\sqrt{-g}\,\mathrm d^4x.
\label{eq:rm-eh-area}
\end{equation}

The scalar response and the tensor response are distinct publications
of the same compliance. The former uses \(4\pi\mathscr C_P\); the latter,
\(8\pi\mathscr C_P\). Upon inserting CT108,

\[
2\pi\mathscr C_P=108\frac{ct_0}{Q},
\qquad
4\pi\mathscr C_P=216\frac{ct_0}{Q},
\qquad
8\pi\mathscr C_P=432\frac{ct_0}{Q}.
\]

The progression \(108\to216\to432\) follows from the geometric factors.
The coincidence of \(432\) with the multiplicity of the transverse region
of \(\pi\) constitutes a structured compatibility; its promotion requires
an incidence map that transports the genealogical multiplicity to the Einstein
tensor.

\section{Ternary Informational Quantum \(\log 3\) and Capacity of the Qutrit Cell}

The thermodynamic relation of the CT108 section is

\begin{equation}
\boxed{Q=k_B\Theta_{\rm clk}\log3.}
\label{eq:rm-planck-trit}
\end{equation}

In this equation, \(\log3\) is the entropy of an equiprobable ternary
alternative; \(\Theta_{\rm clk}\) fixes the thermal chart and \(Q\) publishes the
energy of the cycle. The combination with theorem
\ref{thm:rm-planck-pair} yields

\begin{equation}
\boxed{
Gk_B\Theta_{\rm clk}\log3=c^4\ell.}
\label{eq:rm-gravity-information}
\end{equation}

Gravity transforms the informational quantum of the clock into a length of
geometric response. Equivalently,

\[
\ell=\mathscr C_P k_B\Theta_{\rm clk}\log3.
\]

This equality separates the dimensionless content \(\log3\), the
thermometric frame \(k_B\Theta_{\rm clk}\), and the geometric compliance.

\section{Covariance of bidirectional readers under sheet inversion and temporal conjugation}

Let \(\hbar_{\rm pre}\) and \(\hbar_{\rm ret}\) be the two oriented sections
of the action line, related by a positive quotient
\(R_{\rm act}\):

\[
\frac{\hbar_{\rm pre}}{\hbar_{\rm ret}}=R_{\rm act}.
\]

The invariance of \(\ell\) obliga to

\[
\frac{G_{\rm pre}}{G_{\rm ret}}=R_{\rm act}^{-1}.
\]

Therefore,

\begin{equation}
G_{\sigma}\hbar_{\sigma}=c^3\ell^2
\qquad(\sigma\in\{\mathrm{pre},\mathrm{ret}\}).
\label{eq:rm-two-way-area}
\end{equation}

The action preserves the orientation of the route; its product with the
gravitational response preserves the area. This covariance makes \(\ell^2\)
the geometric invariant of the bidirectional inversion.

\section{Constitutive vacuum and causal closure}

The operator of vacuum proporciona

\[
c_{\rm vac}^{-2}=\mu\varepsilon.
\]

The Planck pair provides a second causal section:

\[
c_{\rm int}=\frac{Q\ell}{\hbar}.
\]

We define the sewing scalar

\begin{equation}
\Xi_c=\frac{Q\ell}{\hbar}\sqrt{\mu\varepsilon}
=\frac{c_{\rm int}}{c_{\rm vac}}.
\label{eq:rm-causal-seam}
\end{equation}

The equality \(\Xi_c=1\) is necessary and sufficient for both readers to
publish the same causal cone. Under this seam,

\begin{equation}
\boxed{
G=\left(\frac{54}{\pi}\right)^2
\frac{t_0^2}{\hbar(\mu\varepsilon)^{5/2}}.}
\label{eq:rm-g-vacuum-product}
\end{equation}

\begin{proof}
The relation \(c=(\mu\varepsilon)^{-1/2}\) gives
\(c^5=(\mu\varepsilon)^{-5/2}\). Substituting it into the circular selector of
\(G\) yields \eqref{eq:rm-g-vacuum-product}.
\end{proof}

The product \(\mu\varepsilon\) governs the causal cone and the gravitational
response; the quotient \(\mu/\varepsilon=Z^2\) governs impedance and
orientation. Both contents follow from the same bilayer operator and preserve
distinct functions.

\section{Machian factorization of the local inertial response through the trace of the global closure: \(\mathcal I_v=\overline{\mathcal I}_v\circ b\)}

A local energy \(E\) determines the dimensionless coordinate

\[
\xi_E(E)=\frac{E}{Q}.
\]

The gravitational coupling of two energies is

\begin{equation}
\alpha_G(E_1,E_2)=\frac{E_1E_2}{Q^2}.
\label{eq:rm-mach-coupling}
\end{equation}

The reference \(Q\) belongs to the global closure clock. Local inertia is therefore evaluated against a scale produced by the complete genealogy. This dependence does not posit an instantaneous action of the universe upon the body: it establishes an operatorial factorization of the local response by the global closure class.

\begin{proofstatusbox}
The identities of the pair, the tension, the compliance, and the covariance
\(G\hbar\) are exact in the declared circular CT108 realization. The
identification \(c_{\rm vac}=c_{\rm int}\) is a typed physical interface;
its content is concentrated in the map of causal lines whose evaluation is
\(\Xi_c=1\).
\end{proofstatusbox}

\begin{falsifierbox}
The closure of this chapter is refuted by: a distinct positive pair that simultaneously satisfies the two governing equations; the loss of \(G\hbar=c^3\ell^2\) under inversion of the action; a section of the vacuum that does not satisfy \(c^{-2}=\mu\varepsilon\); or a frame transport for which \(\Xi_c\ne1\).
\end{falsifierbox}
